% =====================================================================
%  GLOSARIO Y ACRÓNIMOS
%  Uso en el texto:  \gls{clave}   (primera vez y siguientes)
%                    \glspl{clave} (plural)
%  Solo aparecen en el glosario los términos que se usan con \gls.
% =====================================================================

% ---------- Términos del glosario ----------
\newglossaryentry{api}{
	name=API,
	description={Interfaz de Programación de Aplicaciones. Conjunto de reglas y protocolos que permite a diferentes aplicaciones comunicarse entre sí \citapendiente{}}
}

\newglossaryentry{endpoint}{
	name=endpoint,
	description={Punto de acceso de una API o servicio web, identificado por una URL y un método HTTP}
}

\newglossaryentry{token}{
	name=Token,
	description={Credencial digital que el cliente presenta en cada petición para demostrar que está autorizado a acceder a un recurso}
}

% ---------- Acrónimos ----------
\newacronym{tfg}{TFG}{Trabajo de Fin de Grado}
\newacronym{ia}{IA}{Inteligencia Artificial}
\newacronym{iot}{IoT}{Internet of Things}
